% SZKIC: tekst przygotowany z pomocą AI. Przepisz go własnymi słowami, sprawdź źródła i uwzględnij to w deklaracji poziomu AI.
\chapter*{Wstęp}
\addcontentsline{toc}{chapter}{Wstęp}

Płatności elektroniczne należą do podstawowych usług sektora finansowego, a wraz z ich upowszechnieniem rośnie znaczenie wykrywania nadużyć. Każda transakcja musi być oceniona w krótkim czasie, a oszustwa stanowią znikomą część wszystkich operacji. Wykrywanie oszustw płatniczych jest więc jednocześnie problemem klasyfikacji silnie niezbalansowanej, problemem detekcji anomalii oraz problemem inżynierskim: system oceniający transakcje musi działać szybko, być obserwowalny i pozwalać na zarządzanie modelami \cite{dalpozzolo2018}. \todo{1--2 zdania z liczbami o skali oszustw płatniczych, ze źródłem (np. raport EBA/ECB).}

W literaturze porównuje się zwykle metody uczenia nadzorowanego, które korzystają z etykiet oszustw, oraz nienadzorowanego, które szukają transakcji odbiegających od normy bez etykiet \cite{niu2019,carcillo2019}. Wyniki takich porównań zależą od sposobu podziału danych i doboru miar jakości, a sama ocena modeli rzadko jest połączona z narzędziem, w którym użytkownik mógłby zarządzać modelami i obserwować ich działanie.

Celem pracy jest zaprojektowanie i implementacja webowego systemu do wykrywania anomalii w transakcjach płatniczych oraz porównanie skuteczności algorytmów uczenia nadzorowanego i nienadzorowanego na silnie niezbalansowanych danych. Metodykę porównania oparto na pracy \cite{niu2019}, która jest najbliższa tematycznie, i rozszerzono ją o podział czasowy danych, miary właściwe dla klas niezbalansowanych oraz dobór progu na zbiorze walidacyjnym. Wkład własny stanowią: system webowy z modułowym interfejsem do źródeł transakcji, zarządzaniem modelami i monitorowaniem metryk, odtworzenie i rozszerzenie eksperymentów z literatury w ujednoliconej metodyce oraz analiza wpływu metryki, progu i sposobu radzenia sobie z niezbalansowaniem na wnioski z porównania. \todo{doprecyzować po rozmowie z promotorem}

W pracy postawiono następujące pytania badawcze:
\begin{itemize}
  \item na ile modele nadzorowane przewyższają nienadzorowane (PR-AUC, F1) przy podziale czasowym danych,
  \item jak wybór metryki (ROC AUC a PR-AUC) i progu decyzyjnego wpływa na wnioski z porównania,
  \item jak sposób radzenia sobie z niezbalansowaniem klas (wagi klas, SMOTE) wpływa na modele nadzorowane,
  \item czy system działa w czasie zbliżonym do rzeczywistego (opóźnienie oceny pojedynczej transakcji).
\end{itemize}
Hipoteza robocza zakłada, że modele nadzorowane osiągną wyższy PR-AUC, ale nienadzorowane pozostaną użyteczne, zwłaszcza gdy etykiet brakuje lub są opóźnione.

Eksperymenty przeprowadzono na publicznym zbiorze \emph{Credit Card Fraud Detection} \cite{kaggle_ulb}. Praca ma charakter akademicki: \enquote{czas rzeczywisty} jest symulowany przez odtwarzanie zapisanych transakcji, a wdrożenie produkcyjne i integracja z rzeczywistym systemem płatniczym pozostają poza zakresem. Architekturę systemu zaprojektowano modułowo, tak aby możliwe było dołączenie kolejnych źródeł transakcji (własne pliki, interfejsy bankowe), co opisano jako kierunek dalszego rozwoju. \todo{uaktualnić, jeśli PoC z API bankowym wejdzie do zakresu}

W pierwszym rozdziale opisano domenę, czyli operacje płatnicze i problem oszustw. Rozdział drugi przedstawia metody uczenia maszynowego, miary jakości i przegląd literatury. W rozdziałach trzecim i czwartym zaprezentowano projekt oraz implementację systemu, a w piątym badania własne i analizę wyników. Ostatnią częścią pracy jest podsumowanie, w którym zebrano najważniejsze wnioski.
